\documentclass[a4paper,12pt]{article}

\usepackage{fontspec}
\setmainfont{Times New Roman}
\usepackage[russian]{babel}
\usepackage{graphicx, float}
\usepackage{geometry}
\geometry{left=3cm,right=1.5cm,top=2cm,bottom=2cm}
\usepackage{setspace}
\onehalfspacing
\setlength{\parindent}{1.25cm}

\begin{document}

\begin{titlepage}
  \thispagestyle{empty}
  \begin{center}

    
\includegraphics[width=0.18\textwidth]{../assets/gerb.png}

    \vspace{0.5cm}

    {\small
      МИНИСТЕРСТВО НАУКИ И ВЫСШЕГО ОБРАЗОВАНИЯ\\
      РОССИЙСКОЙ ФЕДЕРАЦИИ
    }

    \vspace{0.3cm}

    {\small
      Федеральное государственное бюджетное образовательное учреждение\\
      высшего образования
    }

    \vspace{0.3cm}

    {\large \textbf{«МИРЭА --- Российский технологический университет»}}

    {\normalsize РТУ МИРЭА}

    {\normalsize Институт искусственного интеллекта (ИИИ)}\\
    {\normalsize Кафедра промышленной информатики (ПИ)}

    \vspace{2cm}

    {\LARGE \textbf{ОТЧЁТ}}

    \vspace{0.3cm}

    {\Large \textbf{ПО ПРАКТИЧЕСКОЙ РАБОТЕ № 26}}

    \vspace{0.4cm}

    {\large по дисциплине «Информатика»}

    \vspace{2.5cm}

    \begin{tabular}{ll}
      Выполнил студент группы КВБО-11-25 & Шапаренко Ф.А. \\
      Принял ассистент: & Беднов Г.А. \\
    \end{tabular}

    \vfill

    Практическую работу выполнил: «07» мая 2026 г.\\
    «Зачтено» \rule{1cm}{0.4pt} 2026 г.

    \vspace{1cm}

    Москва 2026

  \end{center}
\end{titlepage}

\noindent\textbf{Тема:} «Разбиение данных, квантование и фильтрация
для  трансформации данных»\\
\noindent\textbf{Цель работы:} \\ ● Изучить методы разбиения данных,
квантования и фильтрации.  \\ ● Применить данные методы для
трансформации и анализа  различных типов данных.

% https://colab.research.google.com/drive/1Dn0CSL72OxtzFeHA0K39orJAYD1gvSlH?usp=sharing#scrollTo=1RWpXLTsDX1E
\begin{figure}[H]
  \centering
  \includegraphics[width=0.5\textwidth]{/home/shafed/projects/02-computer-science/defend-5/assets/26-1.png}
\end{figure}
\begin{figure}[H]
  \centering
  \includegraphics[width=0.8\textwidth]{/home/shafed/projects/02-computer-science/defend-5/assets/26-2.png}
\end{figure}
\begin{figure}[H]
  \centering
  \includegraphics[width=0.8\textwidth]{/home/shafed/projects/02-computer-science/defend-5/assets/26-3.png}
\end{figure}
\begin{figure}[H]
  \centering
  \includegraphics[width=0.8\textwidth]{/home/shafed/projects/02-computer-science/defend-5/assets/26-4.png}
\end{figure}
\begin{figure}[H]
  \centering
  \includegraphics[width=0.8\textwidth]{/home/shafed/projects/02-computer-science/defend-5/assets/26-5.png}
\end{figure}

4.1. Температура повышается от зимы к лету и снижается к декабрю.
Максимум — в июне, минимум — в декабре.

4.2. Квантование упростило анализ, разделив температуры на категории.
Распределение стало нагляднее, но часть точных значений потерялась.

4.3. При фильтрации был обнаружен один выброс. Он отмечен красной
точкой на графике сравнения данных.

4.4. Сглаживание убрало резкие колебания и показало общий тренд:
потепление к лету и похолодание к зиме.

4.5. Обработанные данные проще анализировать, но они менее точные.
Квантование и сглаживание могут скрывать отдельные важные изменения.
\end{document}
